\paragraph{Statement E.16 (finite homology rank in a perfect normal subgroup).}
Let $G$ admit an actual CW classifying space of dimension at most two,
and let $N\triangleleft G$ be perfect. Suppose that
\[
 H_2(N;\mathbb Z)\cong\mathbb Z^r,
 \qquad 0<r<\infty.
\]
Then the quotient $Q=G/N$ acts faithfully on $H_2(N;\mathbb Z)$,
is finite, and satisfies
\[
                         |Q|\mid r+1.
\]
In particular $r=1$ implies $[G:N]\leq2$, and a perfect $G$ then
equals $N$. For every finite positive $r$, if $G$ has no nontrivial
finite image, then $G=N$. At the distinct boundary $r=0$, the
conclusion is instead $\operatorname{cd}_{\mathbb Z}Q\leq2$;
in particular $Q$ is torsion-free. No geometric-dimension assertion
for this quotient is made.

\emph{Proof.}
Let $X$ be the given two-dimensional classifying space and let
$X_N\to X$ be its regular $N$-cover. Its cellular chains are free
$\mathbb ZQ$-modules, without a finite-rank assumption. Since the
cover is a classifying space for $N$, perfection and the absence
of three-cells give the exact sequence
\[
 0\longrightarrow M\longrightarrow C_2(X_N;\mathbb Z)
 \longrightarrow C_1(X_N;\mathbb Z)
 \longrightarrow C_0(X_N;\mathbb Z)
 \longrightarrow\mathbb Z\longrightarrow0,
 \qquad M=H_2(N;\mathbb Z).
 \tag{E.16.1}
\]
This also explains why $H_2(N;\mathbb Z)$ is free abelian.

First $Q$ is torsion. If $q\in Q$ had infinite order, restrict
the rational cellular complex to
$R=\mathbb Q[\langle q\rangle]=\mathbb Q[t,t^{-1}]$.
Its chain modules remain free: restriction from $\mathbb Q Q$
uses coset representatives as a basis. Because there are no
three-cells, $H_2(X_N;\mathbb Q)$ is a submodule of the free
$R$-module $C_2(X_N;\mathbb Q)$. The ring $R$ is a principal
ideal domain, so this submodule is free, including when the
ambient module has infinite rank. A nonzero free $R$-module
has infinite rational dimension, contradicting
$H_2(X_N;\mathbb Q)\cong\mathbb Q^r$.

Next the action of $Q$ on $M$ is faithful. Suppose its kernel
contains an element of prime order $p$, and let $P$ be the
cyclic subgroup it generates. Restrict (E.16.1) to $\mathbb ZP$.
The three chain modules remain free and have zero positive
$P$-cohomology. This is true for arbitrary free ranks: the
cyclic periodic resolution has finite free terms and its
cochain operators $g-1$ and $1+g+\cdots+g^{p-1}$ are exact
in positive degrees on every direct sum of regular modules.
Three applications of the cohomology long exact sequence give
\[
 H^4(P;M)\cong H^1(P;\mathbb Z)=0.
\]
On the other hand the assumed trivial action on $M=\mathbb Z^r$
gives $H^4(P;M)\cong(\mathbb Z/p)^r\ne0$, by the same periodic
resolution. This contradiction excludes every prime-order
element in the kernel. Since $Q$ is torsion, any nonidentity
kernel element would have a prime-order power. Thus
\[
                       Q\hookrightarrow\operatorname{GL}_r(\mathbb Z).
\]

For completeness, a torsion subgroup of
$\operatorname{GL}_r(\mathbb Z)$ is finite. Reduction modulo
$3$ is injective on such a subgroup because its kernel has
no torsion. To check this last assertion, replace a hypothetical
nonidentity torsion matrix in that kernel by a power of prime
order $\ell$ and write it as $A=I+3^sB$, with $s\geq1$ and
$B$ nonzero modulo $3$. If $\ell\ne3$, expansion of $A^\ell$
modulo $3^{s+1}$ contradicts $A^\ell=I$. If $\ell=3$, expansion
modulo $3^{s+2}$ gives the same contradiction: the leading
term is $3^{s+1}B$, while both later terms vanish modulo
$3^{s+2}$. Hence $Q$ injects into the finite group
$\operatorname{GL}_r(\mathbb F_3)$.

It remains to prove the sharper order divisibility. Let $P$ be
a Sylow $p$-subgroup of the now finite group $Q$, put
$k=\mathbb F_p$ and $A=kP$, and take cellular chains over $k$.
The integral universal coefficient sequence gives
\[
 H_1(X_N;k)=0,
 \qquad H_2(X_N;k)=M\otimes k\cong k^r:
\]
the possible degree-two Tor term is
$\operatorname{Tor}_1^{\mathbb Z}(H_1(X_N;\mathbb Z),k)=0$,
and the degree-one term uses the torsion-free group $H_0=\mathbb Z$.
Consequently the restricted $A$-complex gives an exact sequence
\[
 0\longrightarrow M_k\longrightarrow C_2(X_N;k)
 \longrightarrow C_1(X_N;k)\longrightarrow C_0(X_N;k)
 \longrightarrow k\longrightarrow0,
 \qquad \dim_k M_k=r.
 \tag{E.16.2}
\]

Choose a partial resolution of $k$ by finite free $A$-modules
$D_0,D_1,D_2$, and write $L$ for its third syzygy. These
modules and $L$ are finite-dimensional because $A$ is a
finite-dimensional algebra. The iterated Schanuel lemma
applied to (E.16.2) gives
\[
 M_k\oplus D_2\oplus C_1\oplus D_0
 \cong L\oplus C_2\oplus D_1\oplus C_0.
 \tag{E.16.3}
\]
Here the $C_i$ may still have infinite free ranks. The following
finite stabilization argument prevents cancellation of infinite
dimensions from entering the proof.

If finite $A$-modules $U,V$ become isomorphic after addition
of arbitrary free modules, the inclusion and projection maps
of that isomorphism give maps $f:U\to V$, $g:V\to U$ such
that $gf-1_U$ and $fg-1_V$ factor through free modules.
Because $U,V$ are finitely generated, both factorizations
use finite free summands. Write $gf-1_U=ba$ through a finite
free module $F$. Then
\[
 j=(f,a):U\longrightarrow V\oplus F,
 \qquad s=(g,-b):V\oplus F\longrightarrow U
\]
satisfy $sj=1_U$. The complement $K=\ker s$ is finite.
Moreover $1-js$ factors through a finite projective module:
its upper-left block is $1-fg$, and each other block either
has domain or range $F$. Restricting this factorization to
$K$ factors $1_K$ through a finite projective module, so
$K$ is itself projective. We have therefore obtained an
actual finite stabilization
\[
                         V\oplus F\cong U\oplus K.
\]
For a finite $p$-group, $A=kP$ is local and every finite
projective $A$-module is free. One can see locality by
inducting on $|P|$: a central element of order $p$ makes
the ideal it generates minus $1$ nilpotent, and the quotient
is the group algebra of the smaller $p$-group; thus the
augmentation ideal is nilpotent and $A$ has residue field $k$.
Lifting a basis modulo that ideal and applying Nakayama's
lemma proves the asserted freeness of finite projectives.
Thus the finite stabilization shows
\[
                         \dim_k U\equiv\dim_k V\pmod{|P|}.
\]

Apply this to (E.16.3). The finite partial resolution gives
\[
 \dim_k L=\dim_k D_2-\dim_k D_1+\dim_k D_0-1
          \equiv-1\pmod{|P|}.
\]
Hence $r\equiv-1\pmod{|P|}$. Each Sylow subgroup order
divides $r+1$, and their orders are pairwise coprime, proving
$|Q|\mid r+1$.

When $r=1$, this gives $Q=1$ or $Q\cong C_2$; a quotient
of a perfect group is perfect, so a perfect $G$ allows only
the first case. The assertion when $G$ has no nontrivial
finite image follows directly from finiteness of $Q$.
Finally, if $r=0$, (E.16.1) itself is a free
$\mathbb ZQ$-resolution of $\mathbb Z$ of length at most two.
It proves $\operatorname{cd}_{\mathbb Z}Q\leq2$, and the
cyclic periodic resolution excludes nontrivial finite-order
elements in $Q$.\hfill$\square$

This restriction does not settle E2. A countable homology
core need not be normal, and its normal closure need not
retain the finite second-homology rank. In particular, adding
unrestricted auxiliary conjugators and then taking a normal
closure still requires an independent homology calculation.
